\chapter{Conclusion}
\label{ch:conclusion}

Aircraft turnarounds determine whether delay is absorbed on the ground or passed on to the next flight, yet the services they consist of are rarely recorded in a form that permits process analysis. Most airports, however, share a few milestones of every flight through A-CDM, and flight-tracking networks observe when aircraft arrive and leave. This thesis asked whether such sparse, publicly shared observations are enough to calibrate an object-centric simulation of the turnaround, and what the calibrated simulation reveals. It answered this question with a data-driven pipeline and evaluated it on Berlin Brandenburg Airport as a case study.

We linked 1,322 departure records from the BER A-CDM application with OpenSky observations of the same aircraft and enriched 10,212 turnarounds over 41 days with carrier class, aircraft type and observed weather. The linked data revealed how turnarounds transfer delay: a late inbound aircraft is passed on when the scheduled buffer falls below about 45~minutes, each inbound minute adds about 0.43~outbound minutes, departure banks increase runway waiting, and delay travels with the aircraft. A classifier using only information available at landing flags late departure targets with an AUC of about 0.8.

We extended an existing object-centric Petri net of airport ground handling~\cite{ocpn_studio} with carrier-specific service times for cleaning, catering, boarding, deboarding and fueling, observed weather, the local night regulation and schedule-bound departures, and simulated every day of the observed demand once. Carrier-specific timing reproduced the observed target turnaround times much better than the original or carrier-agnostic functions, and a two-parameter calibration reduced the error on three held-out days by 36--72\%. The object-centric analysis showed that the cabin branch determines readiness for most turnarounds, that the baggage hold becomes critical for low-cost turnarounds, and that remote stands add almost 20~minutes in the modeled airport. Process discovery and conformance checking on the generated OCEL logs, together with the Promenade analysis, made these mechanisms and the differences between carrier classes visible. Among the tested scenarios, faster boarding was the most effective lever, additional fuel capacity mainly removed the fuel queue, and moving all low-cost flights to remote stands overloaded the remote apron.

Together, these results show that public milestones, once linked through the aircraft, are sufficient to calibrate an object-centric turnaround simulation at the level of carrier classes, and that object-centric event logs are essential to distinguish local waiting from its effect on the turnaround. Nothing in the pipeline is specific to BER: it applies to any airport that publishes A-CDM milestones and is covered by a flight-tracking network, and more generally to processes whose boundaries are observable while their internal activities are not. For such processes, a calibrated object-centric simulation provides realistic event logs today, and the analyses built on them can be transferred to real logs once these are recorded.

\section{Future Work}

The most direct extension concerns turnarounds under time pressure. In the current model, service times do not depend on how late the aircraft is, so the model transfers too much delay at intermediate buffers. Service times that shorten when the buffer is small would let the model compress turnarounds as real operators do, and the observed delay-transfer curve already provides the calibration target.

A second extension concerns the infrastructure. Replacing the generic stand layout and resource fleets with the airport's own data would turn the relative what-if effects into absolute ones and would allow a direct test of remote-stand operations. Similarly, a longer and continuous A-CDM collection, ideally including arrival and in-block milestones, would allow seasonal calibration, a calibration at the level of individual activities and more precise estimates of weather effects.

Finally, two methodological directions follow from the results. The early-warning classifier could be combined with a simulation of the remaining turnaround, so that predictions are updated at every new milestone during operations. And because the inductive miner represents the highly concurrent handling phase as a flower structure, discovery methods that preserve a known join structure would yield more precise models of strongly concurrent object-centric processes.
